%Reflection{\psset{linestyle=solid}\textcolor{white}{\Large Відбиття}\begin{itemize}	\item Відбиття EMR залежить від довжини хвилі, що демонструється, коли видиме світло та радіохвилі відбиваються від об'єктів, через які рентгенівські промені проходять наскрізь. Мікрохвилі, які мають велику довжину хвилі порівняно з видимим світлом, відбиваються від металевої сітки в мікрохвильовій печі, тоді як видиме світло проходить крізь неї.		%Mirror	\rput(0.5\linewidth,-.6){%Reflection{\psset{linestyle=solid}\textcolor{white}{\Large Відбиття}\begin{itemize}	\item Відбиття EMR залежить від довжини хвилі, що демонструється, коли видиме світло та радіохвилі відбиваються від об'єктів, через які рентгенівські промені проходять наскрізь. Мікрохвилі, які мають велику довжину хвилі порівняно з видимим світлом, відбиваються від металевої сітки в мікрохвильовій печі, тоді як видиме світло проходить крізь неї.		%Mirror	\rput(0.5\linewidth,-.6){%Reflection{\psset{linestyle=solid}\textcolor{white}{\Large Відбиття}\begin{itemize}	\item Відбиття EMR залежить від довжини хвилі, що демонструється, коли видиме світло та радіохвилі відбиваються від об'єктів, через які рентгенівські промені проходять наскрізь. Мікрохвилі, які мають велику довжину хвилі порівняно з видимим світлом, відбиваються від металевої сітки в мікрохвильовій печі, тоді як видиме світло проходить крізь неї.		%Mirror	\rput(0.5\linewidth,-.6){%Reflection{\psset{linestyle=solid}\textcolor{white}{\Large Відбиття}\begin{itemize}	\item Відбиття EMR залежить від довжини хвилі, що демонструється, коли видиме світло та радіохвилі відбиваються від об'єктів, через які рентгенівські промені проходять наскрізь. Мікрохвилі, які мають велику довжину хвилі порівняно з видимим світлом, відбиваються від металевої сітки в мікрохвильовій печі, тоді як видиме світло проходить крізь неї.		%Mirror	\rput(0.5\linewidth,-.6){\input{pictures/reflection.tex}}		\vspace{0.72in}	\item EMR будь-якої довжини хвилі може бути відбите, однак відбивна здатність матеріалу залежить від багатьох факторів, включаючи довжину хвилі падаючого променя.	\item Кут падіння ($\theta_i$) та кут відбиття ($\theta_r$) є однаковими.\end{itemize}}}		\vspace{0.72in}	\item EMR будь-якої довжини хвилі може бути відбите, однак відбивна здатність матеріалу залежить від багатьох факторів, включаючи довжину хвилі падаючого променя.	\item Кут падіння ($\theta_i$) та кут відбиття ($\theta_r$) є однаковими.\end{itemize}}}		\vspace{0.72in}	\item EMR будь-якої довжини хвилі може бути відбите, однак відбивна здатність матеріалу залежить від багатьох факторів, включаючи довжину хвилі падаючого променя.	\item Кут падіння ($\theta_i$) та кут відбиття ($\theta_r$) є однаковими.\end{itemize}}}		\vspace{0.72in}	\item EMR будь-якої довжини хвилі може бути відбите, однак відбивна здатність матеріалу залежить від багатьох факторів, включаючи довжину хвилі падаючого променя.	\item Кут падіння ($\theta_i$) та кут відбиття ($\theta_r$) є однаковими.\end{itemize}}